% Proposed B.147. Owned source; no private PDF/PNG requested.
\subsection*{The complete family link obstructs every injective center extension}

\statement{B.147}{Proved}{The prescribed center basis change has no injective extension to the actual regional kernel, even with arbitrary changes of the family images}
\label{stmt:B147}
Let $N_{\rm reg}=A(\Sigma)$ be B.141's actual ordinary, unfilled
regional kernel. Let $\varphi$ be B.145's automorphism of its
center free subgroup, sending each pivot $Z_n$ to
$u_n=\nu_nZ_n\nu_n^{-1}$ and fixing every complementary center.
There are three specified nontrivial actual family generators
$\lambda_\alpha,\lambda_\delta,\lambda_\epsilon$ such that every
endomorphism $\Phi:N_{\rm reg}\to N_{\rm reg}$ extending
$\varphi$ must satisfy
\begin{equation}\label{eq:B147-forced-killing}
 \Phi(\lambda_\alpha)=\Phi(\lambda_\delta)
                     =\Phi(\lambda_\epsilon)=1.
\end{equation}
Thus no injective endomorphism, and in particular no compatible
basis automorphism, extends the prescribed center change. The
family images may be changed arbitrarily in this assertion;
they are not required to retain their names, conjugacy classes
or abelian coordinates.

Noninjective extensions do exist. This is an ordinary unfilled
obstruction, not a statement about any compact extension, any
different center normalization or any ambient filling collapse.

\begin{proof}
\emph{B.147.1. All family orbits and their exact stabilizers.}
Use B.145's $T_h=F(h,a,d,e,V_4)$ and
$Q_0=\langle h^2,a,d,e,V_4\rangle$. The entire center orbit is
$T_h/Q_0$, with $X_{gQ_0}=gkg^{-1}$. The complete four family
orbits have representatives
$\alpha_1,\beta_0,\delta_1,\epsilon_1$, all adjacent to $k$.
Their stabilizers for the auxiliary $T_h$ action are
\begin{equation}\label{eq:B147-stabilizers}
\begin{aligned}
 Q_\alpha&=\langle hah^{-1},h^2V_4h^{-2},hdh^{-1},heh^{-1}\rangle,\\
 Q_\beta&=\langle a,hV_4h^{-1},hdh^{-1},heh^{-1}\rangle,\\
 Q_\delta=Q_\epsilon&=h\langle a,d,e,V_4\rangle h^{-1}.
\end{aligned}
\end{equation}
Indeed B.141's full retained links are respectively
$(a_1,B_1,d_1,e_1)$, $(a_0,B_0,d_1,e_1)$, and
$(a_1,d_1,e_1,B_0)$ for each of the last two families.
Since $B=hV_4h^{-1}$, these give exactly
\eqref{eq:B147-stabilizers}. A family stabilizer has zero
$h$-height because $h$ shifts its index; the actual center has
period two. Combining all retained translations and all $h$
shifts gives precisely the coset sets $T_h/Q_j$ above. The two
last families are distinct orbits despite their equal stabilizers.

The complete edge rule is that a center coset and a family coset
are adjacent if and only if they have a common representative.
No two centers, and no two families, are adjacent. These facts
are the full B.141 graph-kernel construction, with no height or
word cutoff.

\emph{B.147.2. A complementary center fixed by the prescribed change.}
Retain B.145's exact words
\[
 D=ded^{-1},\quad
 W=ahV_4h^{-1}a^{-1}hDh^{-1}d^{-1},\quad r=WhW^{-1},
\]
\[
 Y_0=X_{WQ_0},\quad Y_1=X_{WhQ_0},\quad
 Z_n=X_{r^nQ_0},\quad q=Y_1^{-1}r.
\]
Put $l=hah^{-1}$ and $X=X_{rlQ_0}$, and define
\[
 \lambda_\alpha=r(\alpha_1),\qquad
 \lambda_\delta=r(\delta_1),\qquad
 \lambda_\epsilon=r(\epsilon_1).
\]
Because $l$ belongs to $Q_\alpha,Q_\delta,Q_\epsilon$, all
three families are adjacent to both $Z_1$ and $X$.

The center $X$ lies outside every pivot coset. Equality
$rlQ_0=r^nQ_0$ would give $r^{1-n}l\in Q_0$. For $n=1$,
$hah^{-1}$ exits the exact two-vertex $Q_0$ core at $a$ from
its nonroot vertex. For $n\ne1$ the word
\[
 Wh^{1-n}W^{-1}hah^{-1}
\]
is reduced at all displayed junctions and retains the initial
$ahV_4$ exit. In the Schreier cover, a reduced path cannot return
from the attached tree. Thus the word is outside $Q_0$ for every
such $n$. It follows that
\begin{equation}\label{eq:B147-center-values}
 \varphi(X)=X,\qquad
 \varphi(Z_1)=Y_1^{-1}Z_1Y_1.
\end{equation}
The centers $X,Z_1,Y_1$ are distinct; the remaining distinctions
also follow from the complete link computation below.

\emph{B.147.3. The entire common link, including every retained coset.}
Let $E$ be the even-$h$ subgroup of $T_h$. Its exact free Schreier
basis is
\[
\begin{gathered}
 H=h^2, \\
 a_0=a,\ a_1=hah^{-1}, \\
 d_0=d,\ d_1=hdh^{-1}, \\
 e_0=e,\ e_1=heh^{-1}, \\
 v_0=V_4,\ v_1=hV_4h^{-1}.
\end{gathered}
\]
Consequently
\begin{equation}\label{eq:B147-free-factors}
\begin{gathered}
 E=Q_0*P, \\
 Q_0=F(H,a_0,d_0,e_0,v_0), \\
 P=F(a_1,d_1,e_1,v_1).
\end{gathered}
\end{equation}
The complete stabilizer types become
\begin{equation}\label{eq:B147-stabilizer-products}
\begin{aligned}
 Q_\alpha&=\langle Hv_0H^{-1}\rangle*F(a_1,d_1,e_1),\\
 Q_\beta&=\langle a_0\rangle*F(v_1,d_1,e_1),\\
 Q_\delta=Q_\epsilon&=P.
\end{aligned}
\end{equation}
Every common family neighbor of $Q_0$ and $lQ_0$ can be written
$gQ_j$ with $g\in Q_0$; its second adjacency is exactly
\begin{equation}\label{eq:B147-coset-condition}
 g^{-1}a_1\in Q_jQ_0.
\end{equation}
Reduced $Q_0*P$ normal form determines all solutions. In the
alpha case the $P$ syllable $a_1$ lies in its indicated free
factor, and the initial $Q_0$ syllable forces
$g\in\langle Hv_0H^{-1}\rangle$. All those $g$ lie in
$Q_\alpha$ and give its single identity coset. In the beta case
its $P$ factor omits the free letter $a_1$, so no solution exists.
For delta and epsilon the $Q_0$ factor is trivial, forcing $g=1$.
Therefore the complete common link is
\begin{equation}\label{eq:B147-complete-link}
 \operatorname{link}(Q_0)\cap\operatorname{link}(lQ_0)
                  =\{\alpha_1,\delta_1,\epsilon_1\}.
\end{equation}
Every $h^{2n}$ translate is included among the representatives
$g\in Q_0$; odd-height family cosets cannot meet $Q_0$ because
all four stabilizers lie in $E$. Thus no additional family level
has been silently omitted.

None of the three families in \eqref{eq:B147-complete-link} is
adjacent to $Y_0=WQ_0$. For alpha, adjacency would require
$W\in Q_\alpha Q_0$. The exponent of $v_1$ vanishes on both
factors, whereas the complete marked word is
\begin{equation}\label{eq:B147-W-normal}
 W=a_0v_1a_0^{-1}(d_1e_1d_1^{-1})d_0^{-1},
\end{equation}
whose $v_1$ exponent is one. For delta and epsilon, adjacency
would require $W\in PQ_0$. Its displayed reduced normal form
has five nontrivial alternating syllables, so this is impossible.
Translation by $r$ now gives
\begin{equation}\label{eq:B147-triple-link}
\begin{gathered}
 \operatorname{link}(X)\cap\operatorname{link}(Z_1)
       =\{\lambda_\alpha,\lambda_\delta,\lambda_\epsilon\},\\
 \operatorname{link}(X)\cap\operatorname{link}(Z_1)
                    \cap\operatorname{link}(Y_1)=\varnothing.
\end{gathered}
\end{equation}
In particular $Y_1$ is distinct from the other two centers.

\emph{B.147.4. The full centralizer intersection.}
For any ordinary graph group $A(\Sigma)$ and three distinct
pairwise nonadjacent vertices $x,z,y$,
\begin{equation}\label{eq:B147-centralizer-lemma}
 C(x)\cap y^{-1}C(z)y
       =A\bigl(\operatorname{link}(x)\cap\operatorname{link}(z)
                                      \cap\operatorname{link}(y)\bigr).
\end{equation}
Here is a normal-form proof valid for the countable graph as well.
Delete $y$ to obtain $G_0=A(\Sigma\setminus\{y\})$.
The whole group is the HNN extension of $G_0$ in which the stable
letter $y$ centralizes the embedded coordinate subgroup
$A(\operatorname{link}(y))$. Coordinate embeddings follow from
their retractions and defining graph presentations. The ordinary
graph reduced-word normal form gives
$C(x)=A(\operatorname{star}(x))$ and the analogous formula for
$z$; both groups lie in $G_0$. If $y^{-1}vy\in G_0$ with
$v\in C(z)$, Britton's lemma forces
$v\in A(\operatorname{link}(y))$, after which the stable letters
cancel. Intersecting the three coordinate subgroups gives
\eqref{eq:B147-centralizer-lemma}; coordinate intersections are
read directly from the same reduced-word normal form.

The centers are pairwise nonadjacent. Applying
\eqref{eq:B147-centralizer-lemma} to
\eqref{eq:B147-triple-link} proves the full group equality
\begin{equation}\label{eq:B147-full-intersection}
 C_{N_{\rm reg}}(X)\cap
 C_{N_{\rm reg}}(Y_1^{-1}Z_1Y_1)=1.
\end{equation}
This is stronger than a nontrivial image in a coordinate quotient:
there is no unexamined complementary family element in this
centralizer intersection.

\emph{B.147.5. Arbitrary family images and exact scope.}
For each $j=\alpha,\delta,\epsilon$, the actual relations
$[X,\lambda_j]=[Z_1,\lambda_j]=1$ hold before filling. If
$\Phi$ extends the prescribed center map, then
\eqref{eq:B147-center-values} puts $\Phi(\lambda_j)$ in the
intersection \eqref{eq:B147-full-intersection}. This proves
\eqref{eq:B147-forced-killing}. Each $\lambda_j$ is a nontrivial
coordinate generator, so $\Phi$ cannot be injective, regardless
of its other family images.

There are nevertheless noninjective extensions: delete every
family generator, apply $\varphi$ on the free center subgroup,
and include that subgroup again. This distinction is necessary.
The B.145 center substitution is still a valid Tietze change in
the complete presentation, with all graph relations rewritten.
The present obstruction does not kill any family element in an
actual filled group, does not evaluate an ambient restriction
kernel, and does not exclude a different center normalization.
No compact graph-group or completed normal-closure inference
has been used.

The complete all-coset proof and independently checked provenance
are retained in
\path{work_DE/next_marginal_source/prescribed_center_no_injective_extension30.md}.
The argument above supplies the full proof, including all integer
and family-coset comparisons, without a translate census.
\end{proof}
